\documentclass[11pt]{article}
\usepackage[margin=1in]{geometry}
\usepackage{booktabs}
\usepackage{graphicx}
\usepackage{float}
\usepackage{hyperref}
\usepackage{caption}
\usepackage{setspace}


\begin{document}
\singlespacing

\section*{Part 4: Preliminary Results and Methods}

\subsection*{Data and Scope}
This prototype analyzes weather--price relationships in PJM. The export has 1,829 records across 49 nodes and 7 regions, covering 2021--2025 (earlyAI/recentAI), with target lmp and weather variables CDH, HDH, TEMP\_C, RH\_PCT, WIND\_MS, DEW\_C, SLP\_HPA, and PRECIP\_MM.

\subsection*{Methods}
Results were generated with the following data-acquisition and processing pipeline:
\begin{enumerate}
    \item \textbf{Energy data source (PJM):} We used the PJM bundle containing day-ahead LMP parquet files (\texttt{pjm\_da\_full\_system\_parquet}), load forecast files, and node metadata (\texttt{feature\_data}). The bundle was downloaded from the shared Box link and placed at \texttt{data/pjm data} (or accessed via \texttt{PJM\_DATA\_ROOT}).
    \item \textbf{Weather data source (NOAA ISD):} For each selected PJM node, we mapped the node to its nearest NOAA ISD station, then downloaded hourly ``global-hourly'' weather CSV files for years 2021--2025. Raw weather files were saved under \texttt{data/noaa\_isd\_global\_hourly/<stationId>/}.
    \item \textbf{Parsing and cleaning:} NOAA coded fields were converted to numeric units: temperature/dew point (tenths of \(^\circ\)C to \(^\circ\)C), sea-level pressure (tenths of hPa to hPa), wind speed (tenths of m/s to m/s), and precipitation (tenths of mm to mm). Sentinel missing values (e.g., 9999/99999 patterns) were set to null. PJM timestamps were parsed as UTC; NOAA timestamps were parsed and localized to UTC.
    \item \textbf{Feature engineering:} We computed \texttt{CDH} and \texttt{HDH} using a base temperature of \(18^\circ\)C and computed relative humidity (\texttt{RH\_PCT}) from temperature and dew point using a Magnus-form approximation.
    \item \textbf{Merge logic:} For each node-year, PJM LMP rows were merged with NOAA hourly rows using nearest-time matching in UTC (\texttt{merge\_asof}) with a 30-minute tolerance. Rows without matched weather values were excluded from correlation calculations.
    \item \textbf{Analysis construction:} For each node-year-variable pair, we computed Pearson correlation with day-ahead LMP, then stored results with identifiers (\texttt{nodeId}, \texttt{regionId}, \texttt{year}, \texttt{period}, \texttt{variable}, \texttt{target}). Period labels were assigned by year rule in code (\texttt{earlyAI} for 2021 and \texttt{recentAI} for 2022+ in this dataset slice).
    \item \textbf{Visualization aggregation:} The two report figures were generated from \texttt{data/exports/correlations\_by\_region\_period.json} by grouping correlations by data-center bucket and variable/year, plotting group means, and adding standard-error bands/bars for uncertainty context.
\end{enumerate}
Key scripts used were \texttt{scripts/export\_pjm\_metadata.py}, \texttt{scripts/noaa\_isd\_pull.py}, and \texttt{scripts/build\_analysis\_exports.py}.

\subsection*{Preliminary Results}
Fig.~\ref{fig:dcbucket} is the primary comparison of correlation shifts (\texttt{recentAI} - \texttt{earlyAI}) for data-center-heavy vs non-heavy nodes.

\begin{figure}[H]
\centering
\includegraphics[width=0.82\textwidth]{heavy_vs_nonheavy_shift.png}
\caption{Mean shift in weather-to-LMP correlation (\texttt{recentAI} - \texttt{earlyAI}) by data-center intensity bucket, with standard-error bars.}
\label{fig:dcbucket}
\end{figure}

\noindent
Figure~\ref{fig:dcbucket} represents how the weather--LMP relationship changed between periods for each bucket. Positive values indicate stronger recent correlation; negative values indicate weaker recent correlation. In our preliminary sample, CDH and HDH increase for both buckets (slightly more in the heavy bucket), while RH\_PCT becomes more negative for both buckets (slightly more in the non-heavy bucket). This indicates measurable temporal shifts, but not a uniform ``heavy regions always weaken more'' effect.

\begin{figure}[H]
\centering
\includegraphics[width=0.82\textwidth]{avg_corr_by_bucket_over_time.png}
\caption{Supplemental trend view: yearly mean weather-to-LMP correlation by bucket (shaded bands are standard error).}
\label{fig:linebucket}
\end{figure}

\noindent
As a supplement, Fig.~\ref{fig:linebucket} shows that both groups follow a similar time pattern (rise to 2022, then decline through 2025), with only small year-to-year gaps between heavy and non-heavy nodes.

\subsection*{Significance to Research Question}
Our research question asks whether data-center-heavy areas exhibit stronger or systematically different weather--price relationships than non-heavy areas. The primary result (Fig.~\ref{fig:dcbucket}) provides partial evidence of difference by variable (e.g., slightly larger CDH/HDH increases in heavy nodes), but not a consistent directional gap across all variables. The supplemental time-series result (Fig.~\ref{fig:linebucket}) reinforces this interpretation: both groups move similarly over time, suggesting common system-wide effects, with only modest group separation. Together, these preliminary findings support the project's central hypothesis that weather sensitivity is changing over time, while also showing that the ``data-center-heavy vs non-heavy'' distinction is nuanced and requires deeper modeling in later sections.

\subsection*{Implementation Access}
\textbf{Code repository:} \url{https://github.com/jaiparmar940/energy-weather-app} \\
\textbf{Data bundle (PJM + features):} \url{https://duke.box.com/s/xtdrmjv1463wf0o60i1rmts0g0urxp01}

\end{document}
